\section{Introduction and the main result}
\label{sec:1}
Energy and momentum are fundamental notions in physics. In lattice field theory
(the best-developed non-perturbative formulation of quantum field theory),
however, the construction of the corresponding Noether current, the
energy--momentum
tensor~\cite{Callan:1970ze,Coleman:1970je,Freedman:1974gs,Joglekar:1975jm}, is
not straightforward~\cite{Caracciolo:1988hc,Caracciolo:1989pt}, because the
translational invariance is explicitly broken by the lattice structure. In a
recent paper~\cite{Suzuki:2013gza}, a possible method to avoid this
complication being inherent in the energy--momentum tensor on the lattice has
been proposed on the basis of the Yang--Mills gradient flow (or the Wilson flow
in the context of lattice gauge
theory)~\cite{Luscher:2010iy,Luscher:2011bx,Luscher:2013cpa}.%
\footnote{In~Refs.~\cite{Giusti:2010bb,Giusti:2012yj,Robaina:2013zmb,%
Giusti:2013sqa,Giusti:2013mxa,Giusti:2014ila}, an interesting method to define
a lattice energy--momentum tensor from shifted boundary conditions has been
developed. In~Ref.~\cite{Suzuki:2012wx}, a method on the basis of the
$\mathcal{N}=1$ supersymmetry has been proposed.} See
Ref.~\cite{Luscher:2013vga} for a recent review on the gradient flow and
Refs.~\cite{Borsanyi:2012zs,Borsanyi:2012zr,Fodor:2012td,Fritzsch:2013je,%
Monahan:2013lwa,Shindler:2013bia,Bonati:2014tqa} for its applications in
lattice gauge theory.

The basic idea of~Ref.~\cite{Suzuki:2013gza} is the following:\footnote{The
following reasoning was inspired by pioneering experimentation by E.~Itou
and~M.~Kitazawa (unpublished).} Consider the pure Yang--Mills theory. The
gradient flow deforms the bare gauge field~$A_\mu(x)$ according to a flow
equation with a flow time~$t$ [Eq.~\eqref{eq:(3.1)} below] and this makes gauge
field configurations ``smooth'' for positive flow times. It can then be shown
to all orders in perturbation theory that any local products of the flowed
gauge field~$B_\mu(t,x)$ for any strictly positive flow time~$t$ is ultraviolet
(UV) finite when expressed in terms of renormalized
parameters~\cite{Luscher:2011bx}. In particular, no multiplicative
renormalization factor is required to make those local products finite. In
other words, they are renormalized composite operators. Such UV-finite
quantities should be ``universal'' in the sense that they are independent of
the UV regularization chosen, in the limit in which the regulator is removed.
This suggests a possibility that by using the gradient flow as an intermediate
tool one may bridge composite operators defined with the dimensional
regularization, with which the translational invariance is
manifest,\footnote{The drawback of the dimensional regularization is, of
course, that it is defined only in perturbation theory.} and those in the
lattice regularization with which one may carry out non-perturbative
calculations.\footnote{We thus implicitly assume that the finiteness of the
flowed fields that can rigorously be proven only in perturbation theory
persists even in the non-perturbative level.} Following this idea, a formula
that expresses the correctly normalized conserved energy--momentum tensor as a
$t\to0$ limit of a certain combination of the flowed gauge field was
derived~\cite{Suzuki:2013gza}. This formula provides a possible method to
compute correlation functions of the energy--momentum tensor by using the
lattice regularization because the universal combination in the formula should
be independent of the regularization. An interesting point is that the small
flow-time behavior of the (universal) coefficients in the formula can be
determined by perturbation theory thanks to the asymptotic freedom. This
implies that if the lattice spacing is fine enough a further non-perturbative
determination of the coefficients is not necessary. (Practically,
non-perturbative determination of those coefficients may be quite useful and
how this determination can be carried out has been investigated
in~Ref.~\cite{DelDebbio:2013zaa}.) Although the validity of the formula
in~Ref.~\cite{Suzuki:2013gza}, especially the restoration of the conservation
law in the continuum limit, still remains to be carefully investigated, the
measurement of the interaction measure (the trace anomaly) and the entropy
density of the Yang--Mills theory at finite temperature on the basis of the
formula~\cite{Asakawa:2013laa} shows encouraging results; the method appears to
be promising even practically.

In~Ref.~\cite{Suzuki:2013gza}, the method was developed only for the pure
Yang--Mills theory. It is then natural to ask for wider application if the
method can be generalized to gauge theories containing matter fields,
especially fermion fields. In the present paper, we work out this
generalization. We thus suppose a vector-like gauge theory\footnote{We assume
that the theory is asymptotically free.} with a gauge group~$G$ that contains
\begin{equation}
   \text{$N_{\mathrm{f}}$ Dirac fermions in the gauge representation~$R$}.
\label{eq:(1.1)}
\end{equation}
For simplicity, we assume that all $N_{\mathrm{f}}$ fermions possess a common mass; this
restriction might be appropriately relaxed.

Our main result is Eq.~\eqref{eq:(4.70)}, and a step-by-step derivation of this
master formula is given in subsequent sections. For those who are interested
mainly in the final result, here we give a brief explanation of how to read the
master formula~\eqref{eq:(4.70)}: The left-hand side is the
correctly normalized conserved energy--momentum tensor (with the vacuum
expectation value subtracted); our formula~\eqref{eq:(4.70)} holds only when
the energy--momentum tensor is separated from other operators in correlation
functions in the position space. The combinations in the right-hand side are
defined by~Eqs.~\eqref{eq:(4.1)}--\eqref{eq:(4.5)}. There, $G_{\mu\nu}^a(t,x)$
and~$D_\mu$ are the field strength and the covariant derivative of the flowed
gauge field, respectively (the definition of the flowed gauge field is
identical to that
of~Refs.~\cite{Luscher:2010iy,Luscher:2011bx,Luscher:2013cpa}); our
\emph{ringed\/} flowed fermion fields
$\mathring{{\chi}}(t,x)$ and~$\mathring{\Bar{\chi}}(t,x)$ are, on the other
hand, somewhat different from those of~Ref.~\cite{Luscher:2013cpa}, $\chi(t,x)$
and~$\Bar{\chi}(t,x)$, and they are related by~Eqs.~\eqref{eq:(3.20)}
and~\eqref{eq:(3.21)}. The coefficient functions $c_i(t)$
in~Eq.~\eqref{eq:(4.70)} are given
by~Eqs.~\eqref{eq:(4.72)}--\eqref{eq:(4.76)}. There, $\Bar{g}(q)$ and
$\Bar{m}(q)$ are the running coupling and the running mass parameter defined by
Eqs.~\eqref{eq:(4.20)} and~\eqref{eq:(4.21)}, respectively; throughout this
paper, $m$ denotes the (common) renormalized mass of the fermions.
Equations~\eqref{eq:(4.72)}--\eqref{eq:(4.76)} are for the minimal subtraction
($\text{MS}$) scheme and the result for the modified minimal subtraction
($\overline{\text{MS}}$) scheme can be obtained by the
replacement~\eqref{eq:(4.77)}. $b_0$ and~$d_0$ are the first coefficients of
renormalization group functions, Eq.~\eqref{eq:(2.16)}
and~Eq.~\eqref{eq:(2.18)}, respectively. The energy--momentum tensor is given
by the $t\to0$ limit in the right-hand side of~Eq.~\eqref{eq:(4.70)}. As in the
pure Yang--Mills case mentioned above, the combination in the right-hand side
is UV finite and one may use the lattice regularization to compute the
correlation functions of the combination in the right-hand side. In this way,
correlation functions of the correctly normalized conserved energy--momentum
tensor are obtained. To ensure the ``universality,'' however, the continuum
limit has to be taken before the $t\to0$ limit. Practically, with a finite
lattice spacing~$a$, the flow time~$t$ cannot be taken arbitrarily small to
keep the contact with the continuum physics. Instead, we have a natural
constraint,
\begin{equation}
   a\ll\sqrt{8t}\ll R,
\label{eq:(1.2)}
\end{equation}
where~$R$ denotes a typical physical scale, such as the hadronic scale or the
box size. The extrapolation for~$t\to0$ thus generally requires a sufficiently
fine lattice.

Here is our definition of the quadratic Casimir operators: We set the
normalization of anti-Hermitian generators~$T^a$ of the representation~$R$
as~$\tr_R(T^aT^b)=-T(R)\delta^{ab}$ and~$T^aT^a=-C_2(R)1$. We also denote
$\tr_R(1)=\dim(R)$. From the structure constants in~$[T^a,T^b]=f^{abc}T^c$, we
define $f^{acd}f^{bcd}=C_2(G)\delta^{ab}$. For example, for the fundamental
$N$~representation of~$SU(N)$ for which $\dim(N)=N$, the conventional
normalization is
\begin{equation}
   C_2(SU(N))=N,\qquad T(N)=\frac{1}{2},\qquad
   C_2(N)=\frac{N^2-1}{2N}.
\label{eq:(1.3)}
\end{equation}
